\section{Conclusion and future work}
\label{sec:conclusion}

Dynamic authentication graphs formulated strictly over host-to-host edges omit the client
execution context available at Zero Trust policy decision points. Decomposing requests into a
5-node relational chain with an explicit client-configuration node enables detection of
credential reuse from unfamiliar tools as a structural anomaly.

Two broader deployment principles emerge from the streaming evaluation. First, streaming
memory updates must be decoupled from the detector's own predictions through an external
commit gate to prevent runaway baseline drift and adversarial poisoning. Second, synthetic
security benchmarks require systematic leakage audits to prevent models from exploiting
unconstrained synthetic artifacts. On the verified benchmark, the 5-node TGN exceeds
the single-feature baseline floor on lateral movement under one-class training.

\ifpreliminary
The immediate priority is completing reproduction runs: every quantitative claim in this draft
predates generator de-leakage and must be reproduced on the verified generator before release.
\else
Future work spans three operational directions.
\fi
First, evaluating the architecture across multi-source enterprise testbeds such as the AIT
corpus will test cross-environment generalization with packet-derived TLS fingerprints.
Second, evaluating bipartite degradation on traditional benchmarks like LANL under standardized
community splits provides a direct baseline for the incremental utility of fine-grained contextual
nodes. Finally, evaluating distributed state-eviction schemes for millions of concurrent entities
remains an immediate step for enterprise-scale serving.
